% Source: Lee, "Introduction to Riemannian Manifolds" (2nd ed.), Chapter 1 "What Is Curvature?", PDF pages 15-22 of the cited source.
% Transcribed from the cited source, with manual checks for a few pages, and organized according to leanblueprint
% conventions (semantic \label, bracketed titles, \uses dependency edges, \dcref to Lee's own numbering) below.

\section{The Euclidean Plane}

\begin{definition}[rigid motion, congruence]
\label{def:rigid-motion}
\lean{AffineIsometryEquiv}

\mathlibok
A \emph{rigid motion} of a Euclidean space $E$ is a bijection $f\colon E\to E$ such that
\[
\norm{f(x)-f(y)}=\norm{x-y}\qquad (x,y\in E).
\]
Families $(v_i)_{i\in I}$ and $(w_i)_{i\in I}$ of points are \emph{congruent} when a rigid motion
$f$ satisfies $f(v_i)=w_i$ for every $i$; subsets are congruent when a rigid motion maps one onto
the other. By the Mazur--Ulam theorem every distance-preserving bijection of $E$ is affine, so rigid
motions are precisely affine isometric automorphisms.
\end{definition}

\begin{lemma}[isometry extension]
\label{lem:isometry-extension}
\lean{LeeLib.Ch01.exists_affineIsometryEquiv_of_dist_eq}
\uses{def:rigid-motion}
\leanok



Let $E$ be finite-dimensional and Euclidean, and let $(v_i)_{i\in I}$ and $(w_i)_{i\in I}$ be
nonempty point families satisfying
\[
\norm{v_i-v_j}=\norm{w_i-w_j}\qquad (i,j\in I).
\]
Then the two families are congruent.
\end{lemma}
\begin{proof}
\uses{def:rigid-motion}
\sketch Choose $i_0\in I$ and set $v_i'=v_i-v_{i_0}$ and $w_i'=w_i-w_{i_0}$. Polarization gives
\[
2\inner{v_i'}{v_j'}=\norm{v_i'}^2+\norm{v_j'}^2-\norm{v_i'-v_j'}^2,
\]
and the distance hypothesis gives the same expression for the $w_i'$, hence
$\inner{v_i'}{v_j'}=\inner{w_i'}{w_j'}$. Thus the map
\[
\sum_i c_i v_i'\longmapsto\sum_i c_i w_i'
\]
is a well-defined linear isometry from $S=\operatorname{span}\{v_i'\}$ into $E$: equality of the
two squared norms shows that the map is independent of the chosen finite representation. Extend it
to a linear isometry $g\colon E\to E$ by choosing orthonormal complements; finite dimensionality
makes $g$ bijective. Then $f(x)=g(x-v_{i_0})+w_{i_0}$ is a rigid motion with $f(v_i)=w_i$.
\end{proof}

\begin{theorem}[Side-Side-Side]
\label{thm:side-side-side}
\lean{LeeLib.Ch01.side_side_side_iff}
\uses{def:rigid-motion}
\leanok
\dcref{ch1:1.1}



Two Euclidean triangles are congruent if and only if their corresponding side lengths agree.
\end{theorem}
\begin{proof}
\uses{lem:isometry-extension,def:rigid-motion}
\sketch A rigid motion preserves all pairwise distances, proving the forward implication. For the
converse, label the vertices $(a,b,c)$ and $(a',b',c')$. Equality of the three corresponding side
lengths, together with symmetry and vanishing on the diagonal, gives equality of all pairwise
distances for the two indexed triples. Apply Lemma \ref{lem:isometry-extension}.
\end{proof}

\begin{theorem}[Angle-Sum Theorem]
\label{thm:angle-sum-theorem}
\lean{EuclideanGeometry.angle_add_angle_add_angle_eq_pi}
\dcref{ch1:1.2}

\mathlibok
The interior angles of a Euclidean triangle sum to $\pi$.
\end{theorem}

\begin{theorem}[Circle Classification Theorem]
\label{thm:circle-classification}
\lean{LeeLib.Ch01.euclidean_circle_congruent_iff_radius_eq}
\uses{def:rigid-motion}
\leanok
\dcref{ch1:1.3}



Two Euclidean circles are congruent if and only if their radii are equal.
\end{theorem}
\begin{proof}
\uses{def:rigid-motion}
\sketch A circle of radius $R\geq0$ has diameter $2R$: antipodal points realize this distance, while the
triangle inequality bounds every pair of points by $2R$. A rigid motion preserves diameters, so
congruent circles have equal radii. If the radii agree, translation by the difference of the centers
maps one circle onto the other and is a rigid motion.
\end{proof}

\begin{theorem}[Circumference Theorem]
\label{thm:circumference-theorem}
\lean{LeeLib.Ch01.circumference_circle}
\leanok
\dcref{ch1:1.4}


The circumference of a Euclidean circle of radius $R$ is $2\pi R$.
\end{theorem}
\begin{proof}
\sketch For $R\geq0$, parametrize the circle by $\gamma(t)=c+Re^{it}$ for $0\leq t\leq2\pi$ and define its
length as the supremum of polygonal subdivision sums. Since
\[
\norm{\gamma(x)-\gamma(y)}=2R\left|\sin\frac{x-y}{2}\right|\leq R|x-y|,
\]
every subdivision sum is at most $2\pi R$. For the regular subdivision $t_k=2\pi k/n$, the sum
is
\[
2nR\sin\frac{\pi}{n}=2\pi R\frac{\sin(\pi/n)}{\pi/n},
\]
which tends to $2\pi R$; the supremum is therefore exactly $2\pi R$.
\end{proof}

A \emph{curve} is a smooth parametrized map of a real variable; a prime denotes differentiation in
the parameter. For a unit-speed plane curve $\gamma$, its (unsigned) curvature is
\[
\kappa(t)=\norm{\gamma''(t)}.
\]
At $p=\gamma(t)$, tangent circles have centers on the line through $p$ orthogonal to $\gamma'(t)$.
There is a unique unit-speed tangent circle whose acceleration at $p$ equals $\gamma''(t)$, the
\emph{osculating circle}; if $\gamma''(t)=0$, it is replaced by a straight line (infinite radius).
Thus $\kappa=1/R$ for osculating radius $R$; a radius-$R$ circle has curvature $1/R$, while a
straight line has curvature zero. Given a continuous unit normal field $N$, the \emph{signed curvature}
$\kappa_N$ is positive when the curve turns toward $N$ and negative when it turns away.

\begin{theorem}[Plane Curve Classification Theorem]
\label{thm:plane-curve-classification}
\notready
\dcref{ch1:1.5}
\notready
Suppose $\gamma,\widetilde\gamma\colon[a,b]\to\mathbb R^2$ are smooth unit-speed curves with unit
normal fields $N,\widetilde N$. They are congruent by a direction-preserving congruence if and only if
\[
\kappa_N(t)=\kappa_{\widetilde N}(t)\qquad(t\in[a,b]).
\]
\end{theorem}

\begin{theorem}[Total Curvature Theorem]
\label{thm:total-curvature}
\notready
\dcref{ch1:1.6}
\notready
If $\gamma\colon[a,b]\to\mathbb R^2$ is a unit-speed simple closed curve with
$\gamma'(a)=\gamma'(b)$ and $N$ is the inward-pointing normal, then
\[
\int_a^b\kappa_N(t)\,dt=2\pi.
\]
\end{theorem}

\section{Surfaces in Space}

For a surface $S\subseteq\mathbb R^3$, a point $p\in S$, and a unit normal $N$ at $p$, a normal
section is obtained by intersecting $S$ near $p$ with a plane through $p$ parallel to $N$. The signed
curvature of this plane curve, computed with respect to $N$, varies with the normal plane; its minimum
and maximum are the \emph{principal curvatures} $\kappa_1$ and $\kappa_2$ at $p$.

A property of surfaces in $\mathbb R^3$ is \emph{intrinsic} when it is preserved by isometries,
that is, maps preserving lengths of curves. Principal curvatures are not intrinsic: for
\[
S_1=\{(x,y,0):0<x<\pi,\ 0<y<\pi\},\qquad
S_2=\{(x,y,z):z=\sqrt{1-y^2},\ 0<x<\pi,\ |y|<1\},
\]
the principal curvatures (for a downward normal) are $(0,0)$ and $(0,1)$, respectively, but
$(x,y,0)\mapsto(x,\cos y,\sin y)$ is an isometry $S_1\to S_2$.

The Theorema Egregium asserts that
\[
K=\kappa_1\kappa_2
\]
is intrinsic; see \cite{Gau65} and \cite[Vol.~2]{Spi79}. The plane and the half-cylinder above have
$K=0$ (for the latter, $0\cdot1=0$), while a sphere of radius $R$ has principal curvatures
$\pm1/R$ and Gaussian curvature $K=1/R^2$, independent of the choice of unit normal. The constant
Gaussian-curvature model surfaces include the Euclidean plane ($K=0$), spheres of radius $R$
($K=1/R^2$), and the hyperbolic plane ($K<0$), which has no global smooth embedding in $\mathbb R^3$
[Spi79, Vol.~3, pp.~373--385].

\begin{theorem}[Uniformization Theorem]
\label{thm:uniformization-theorem}
\notready
\dcref{ch1:1.7}
\notready
Every connected $2$-manifold is diffeomorphic to a quotient of a constant-Gaussian-curvature model
surface by a discrete fixed-point-free group of isometries. Consequently, every connected
$2$-manifold admits a complete Riemannian metric of constant Gaussian curvature.
\end{theorem}

\begin{theorem}[Gauss--Bonnet Theorem]
\label{thm:gauss-bonnet-theorem-preview}
\notready
\dcref{ch1:1.8}
\notready
If $S$ is a compact Riemannian $2$-manifold, then
\[
\int_S K\,dA=2\pi\chi(S),
\]
where $dA$ is the Riemannian area measure and $\chi(S)$ is the Euler characteristic.
\end{theorem}

The Gauss--Bonnet and uniformization theorems imply that a compact connected orientable surface
admitting a metric with strictly positive Gaussian curvature is diffeomorphic to the sphere. A compact
connected orientable surface with nonpositive Gaussian curvature is not a sphere, and its universal
cover is topologically the plane.

\section{Curvature in Higher Dimensions}

Let $(M,g)$ be an $n$-dimensional Riemannian manifold. Geodesics, locally shortest curves between
nearby points, replace plane sections as intrinsic probes of curvature. Through each $p\in M$ and
$v\in T_pM$ there is a unique geodesic with initial point $p$ and initial velocity $v$.

For $p\in M$ and a $2$-plane $\Pi\subset T_pM$, the geodesics through $p$ with initial velocities
in $\Pi$ locally sweep out a $2$-submanifold $S_\Pi$ with the induced metric. Its Gaussian curvature
at $p$ is the \emph{sectional curvature} $\sec(\Pi)$, so curvature at $p$ is the map
\[
\sec\colon\{\text{2-planes in }T_pM\}\to\mathbb R.
\]

The constant-sectional-curvature model spaces are Euclidean space ($\sec=0$), the sphere of radius
$R$ ($\sec=1/R^2$), and hyperbolic space of radius $R$ ($\sec=-1/R^2$). Not every manifold admits a
constant-sectional-curvature metric.

\begin{theorem}[Characterization of Constant-Curvature Metrics]
\label{thm:characterization-constant-curvature}
\notready
\dcref{ch1:1.9}
\notready
The complete connected $n$-dimensional Riemannian manifolds of constant sectional curvature are,
up to isometry, exactly the quotients $\widetilde M/\Gamma$, where $\widetilde M$ is Euclidean,
spherical, or hyperbolic space with constant sectional curvature and $\Gamma$ is a discrete group of
isometries acting freely on $\widetilde M$.
\end{theorem}

\begin{theorem}[Cartan--Hadamard]
\label{thm:cartan-hadamard-preview}
\notready
\dcref{ch1:1.10}
\notready
If $M$ is a complete connected Riemannian $n$-manifold with sectional curvature everywhere at most
zero, then its universal covering space is diffeomorphic to $\mathbb R^n$.
\end{theorem}

\begin{theorem}[Myers]
\label{thm:myers}
\notready
\dcref{ch1:1.11}
\notready
If $M$ is a complete connected Riemannian manifold whose sectional curvature is bounded below by a
positive constant, then $M$ is compact and has finite fundamental group.
\end{theorem}
